\chapter{Pendahuluan}
Pendahuluan menjelaskan apa yang akan dikerjakan dalam penelitian dan mengapa
penelitian tersebut perlu dikerjakan.

\section{Latar Belakang}
Uraikan latar belakang munculnya ide penelitian. Mulailah dari konteks umum,
lalu persempit menuju masalah khusus yang akan diselesaikan (teknik piramida
terbalik). Dukung setiap klaim dengan fakta dari pustaka atau pengamatan,
misalnya hasil penelitian terdahulu \cite{lecun1998gradient,vaswani2017attention}.

Akhiri bagian ini dengan pernyataan yang menjembatani latar belakang dengan
rumusan masalah. Istilah asing yang belum diserap ditulis miring, misalnya
\textit{machine learning}.

\section{Rumusan Masalah}
Tuliskan pertanyaan penelitian dalam kalimat tanya. Pertanyaan yang baik
memiliki ciri sebagai berikut:
\begin{enumerate}
  \item Jelas: memakai bahasa Indonesia baku dan mudah dipahami.
  \item Relevan: sesuai dengan hal yang diteliti dan konteks keilmuannya.
  \item Fokus: terarah pada masalah yang ingin diselesaikan.
  \item Dapat terjawab: dapat diukur melalui penelitian sesuai batas waktu dan
        sumber daya.
\end{enumerate}

\section{Tujuan Penelitian}
Tuliskan tujuan dalam kalimat pernyataan yang sederhana, jelas, dan selaras
dengan setiap rumusan masalah.

\section{Manfaat Penelitian}
Uraikan dampak positif penelitian bagi ruang lingkup yang lebih luas atau bagi
pemangku kepentingan. Jangan menuliskan ``untuk memenuhi syarat memperoleh gelar
sarjana'' sebagai manfaat.

\section{Batasan Masalah}
Nyatakan batasan dan asumsi yang menjelaskan ruang lingkup masalah. Batasan yang
sangat teknis lebih tepat diletakkan di bab metodologi.
